\section{Introduction}

Reverse docking and multi-target virtual screening turn many separate docking
calculations into one ligand-by-target matrix
\cite{luo2017,lapillo2019,lee2020crds,krause2023}. These matrices are now large.
DOCKSTRING reports AutoDock Vina scores for 260,155 ligands against 58 targets
\cite{vina2010,dockstring2022}, and a published secondary-pharmacology panel
scores 12,651 compounds against 44 receptors \cite{nikitin2025}. Once a matrix
that size exists, people compare its columns. A correlation between two target
columns gets read as a statement about which receptors a chemical series will hit
together, and such correlations are used to group receptors, or to place a
compound in a profile space defined by many pockets at once
\cite{fukunishi2005,fukunishi2006,lee2012profiles,mangione2019cando}. Read that
way, the matrix stops being a stack of independent predictions and becomes a map
of relationships among targets.

That map inherits a known bias. Empirical scoring functions reward size and
contact count, and receptors differ in how much room their pockets have and in
the range of scores they return \cite{pan2003,vigers2004,jacobsson2006}. A large
compound therefore scores well against most pockets, and much of the covariance
between target columns is a single component shared across the whole panel. In
the two matrices we examine, the leading principal component of the target
correlation matrix carries 73.3\% of the variance on the 44-receptor panel and
65.9\% on the 58-target DOCKSTRING panel, and in both cases its loading vector
points almost straight at the direction in which all targets move together:
cosine 0.970 and 0.969 to that uniform direction. Two receptors can look
correlated simply because both reward bulk.

Corrections for this are old. Per-target normalisation, and two-way schemes that
also normalise on the ligand side, have been in use since the early 2000s
\cite{pan2003,carta2007,sepresa2009,casey2009,luo2017}, and principal components
of docking profiles have served as compound and protein descriptors for just as
long \cite{fukunishi2005,fukunishi2006,fukunishi2018}. We use the simplest member
of that family: subtract each ligand's mean across targets, so that what is left
says how a target treats a compound compared with how the rest of the panel
treats it. To summarise a map we report participation-ratio dimension
\cite{mazzucato2016}, which counts how many independent directions the
correlations spread over, and is a re-expression of the mean squared correlation
between target pairs. Neither step is new. Our contribution is the evidence needed
to interpret the resulting map.

Three gaps remain. Single-program studies do not establish whether the shared component
generalises across scoring-function lineages. Removing a per-ligand mean also raises dimension
by construction, so any residual structure must be compared with null matrices that undergo the
same transformation. Finally, a corrected map is useful only if it contributes information beyond
the sequence identity and family labels that practitioners already use. A practical question sits
alongside these: how many ligands must be docked before the estimated map settles?

We address these gaps with a fixed-pose scorer comparison, four transformation-matched nulls,
chemical-domain controls, external pharmacology and sequence baselines, and a sampling-cost analysis.
The shared axis appears across the tested scorers. The centred map retains structured geometry, but
that geometry changes with the ligand library and its external agreement varies by panel
(Figures~\ref{fig:axis}--\ref{fig:external}). A few hundred ligands from the source-library domain
closely recover the full maps (Figure~\ref{fig:cost}). We therefore treat the residual map as a
library-conditional diagnostic, not as an intrinsic protein network or a compound-level target-ranking
rule.
